% AY2020 力学 第1問 転記（問題のみ）
\section*{第1問〔I〕}

図1のように，なめらかな床の上で，ばね定数 $k$ のばねに繋がれた質量 $m$ の物体が，
速度に比例する空気抵抗と外力 $F\cos\omega t$ を受けて往復運動している。空気抵抗の
比例定数を $B\,(>0)$ とし，ばねが自然長のときの物体の重心位置を $x=0$ とする。
このとき，次の問い (1)〜(3) に答えよ。

\begin{enumerate}[label=(\arabic*)]
  \item 物体に関する運動方程式を求めよ。
  \item 物体に関する運動方程式の特解を求めよ。
  \item 物体が運動を始めて，十分時間が経過した後に物体の往復運動の振幅が最大となる
        ときの $\omega$ を求めよ。
\end{enumerate}

\begin{center}
% 図1：なめらかな床の上、ばねにつながれた質量m。x=0が自然長
\begin{tikzpicture}[>=stealth, scale=1.0]
  \draw (-0.2,0) -- (5.0,0);                      % 床
  \fill[pattern=north east lines] (-0.45,0) rectangle (0,1.3);
  \draw (0,0) -- (0,1.3);                          % 壁
  \draw[decorate,decoration={coil,segment length=5pt,amplitude=4pt}]
        (0,0.55) -- (2.0,0.55);                     % ばね
  \draw (2.0,0.1) rectangle (3.2,1.0);             % 物体
  \draw[dashed] (2.6,0.1) -- (2.6,-0.45);
  \node[below] at (2.6,-0.45) {$x=0$};
  \node at (4.5,-0.3) {図 1};
\end{tikzpicture}
\end{center}
